\documentclass[UTF8]{ctexart}

\usepackage{amscd,amsthm,amsfonts,amssymb,amsmath}
\usepackage[colorlinks=true]{hyperref}
\usepackage[all]{xy}
\usepackage{mathrsfs}
\usepackage{tikz}
\usepackage{bm}
\usepackage{geometry}
\usepackage{xcolor}
\usepackage{eso-pic}
% \usepackage{CJK}
\usepackage{arydshln}
% \usepackage[11pt]{moresize}



\newtheorem{thm}{定理}
\newtheorem{cor}[thm]{推论}
\newtheorem{lem}[thm]{引理}
\newtheorem{prop}[thm]{命题}
\newtheorem{defn}[thm]{定义}
\newtheorem{ques}[thm]{问题}
\newtheorem{eg}[thm]{例题}
\newtheorem{ex}[thm]{习题}
\newtheorem{rmk}[thm]{注释}
\newtheorem{ntn}[thm]{记号}




\begin{document}


\title{习题课材料（五）}
\date{}

\maketitle

%\noindent{\Large\textbf{注: 带 $\heartsuit $号的习题有一定的难度、比较耗时, 请量力为之.}}

\begin{ex}
  给定 \(\boldsymbol{a} \in \mathbb{R}^3\).
  计算 \(\boldsymbol{a}\) 与坐标向量 \(\boldsymbol{e}_1,\boldsymbol{e}_2,\boldsymbol{e}_3\)
  的夹角的余弦, 并计算这三个余弦值的平方和.
\end{ex}

\begin{proof}
  设 \(\boldsymbol{a}=[a_1,a_2,a_3]^{\mathrm{T}}\).
  回忆 \(\boldsymbol{a}\) 与 \(\boldsymbol{e}_i\) 的夹角的余弦为
  \(\frac{\boldsymbol{a}^{\mathrm{T}}\boldsymbol{e}_i}{\|\boldsymbol{a}\|\|\boldsymbol{e}_i\|}=\frac{a_i}{\sqrt{a_1^2+a_2^2+a_3^2}}\).
  它们的平方和因此为 \(1\).
\end{proof}

\begin{ex}
  证明, 分块上三角矩阵
  \(X=\begin{bmatrix} c & \boldsymbol{a}^{\mathrm{T}} \\ \boldsymbol{0} & Q\end{bmatrix}\)
  是正交矩阵时, 必有 \(c = \pm 1\), \(\boldsymbol{a} = \boldsymbol{0}\),
  \(Q\) 是正交矩阵.
\end{ex}

\begin{proof}
  进行分块矩阵的乘法运算:
  \begin{align*}
    & X^{\mathrm{T}} \cdot X \\
    =& \begin{bmatrix} c & \boldsymbol{0}^{\mathrm{T}} \\ \boldsymbol{a} & Q^{\mathrm{T}}\end{bmatrix} \cdot
       \begin{bmatrix} c & \boldsymbol{a}^{\mathrm{T}} \\ \boldsymbol{0} & Q\end{bmatrix}  \\
    =&
       \begin{bmatrix}
         c^2 & c\boldsymbol{a}^{\mathrm{T}} \\ c\boldsymbol{a} & Q^{\mathrm{T}} \cdot Q
       \end{bmatrix}
  \end{align*}
  若 \(X\) 是正交矩阵, 则有 \(c^2 = 1\), \(c \cdot \boldsymbol{a} = \boldsymbol{0}\),
  \(Q^{\mathrm{T}}\cdot Q = \mathrm{I}\). 故而问题得证.
\end{proof}

\begin{ex}[\(\heartsuit\)]
  设可逆矩阵 \(A\) 保持向量之间角度不变.
  证明 \(A\) 是正交矩阵的常数倍.
\end{ex}

\begin{proof}
  设 \(A = QR\) 为 \(A\) 的 QR 分解, 其中 \(Q\) 为正交矩阵, \(R\) 为上三角矩阵.
  容易验证, \(R=Q^{-1}A\)仍然为保角变换. 由于
  \(\boldsymbol{e}_1,\ldots,\boldsymbol{e}_n\)
  两两正交, \(R\boldsymbol{e}_1,\ldots,R\boldsymbol{e}_n\) 也两两正交.
  从而 \(R\) 是列正交矩阵, 即 \(R^{\mathrm{T}}R\) 是对角矩阵.
  利用归纳法, 可证明 \(R\) 是对角矩阵.
  设 \(R = \mathrm{diag}(\lambda_1,\ldots,\lambda_n)\).
  若 \(\lambda_1 \neq \lambda_2\),
  则 \(R(x\boldsymbol{e}_1+y\boldsymbol{e}_2)\) 与 \(R\boldsymbol{e}_1\)
  的夹角的余弦为
  \begin{equation*}
    \frac{\lambda_1 x}{\sqrt{\lambda_1^2 x^2 + \lambda_2^2 y^2}},
  \end{equation*}
  而 \(xe_1 + ye_2\) 与 \(e_1\) 的夹角余弦为
  \begin{equation*}
    \frac{x}{\sqrt{x^2 + y^2}}.
  \end{equation*}
  若二者绝对值对一切 \(x,y\) 都相等, 只能有 \(\lambda_1 = \pm\lambda_2\).
  从而 \(R = k\mathrm{diag}(\pm1,\ldots,\pm 1)\) 是数值矩阵乘以某个对角为正负一
  的矩阵. 从而 \(A\) 为数值矩阵乘以某个正交矩阵.
\end{proof}

\begin{ex}
 设 \(f\colon \mathbb{R}^n \to \mathbb{R}\)
  为线性映射. 证明存在唯一 \(\boldsymbol{a} \in \mathbb{R}^n\),
  使得 \(f(\boldsymbol{x}) = \boldsymbol{a}^{\mathrm{T}}\boldsymbol{x}\).
\end{ex}

\begin{proof}
  设 \(f(\boldsymbol{e}_i) = a_i\).
  则 \(f(x_1 \boldsymbol{e}_1 + \cdots + x_n \boldsymbol{e}_n)=a_1 x_1 + \cdots + a_n x_n\).
  即
  \begin{equation*}
    f
    \begin{bmatrix}
      x_1 \\ \vdots \\ x_n
    \end{bmatrix}
    = [a_1 \; \cdots\; a_n]  \cdot
    \begin{bmatrix}
      x_1 \\ \vdots \\ x_n
    \end{bmatrix}. 
  \end{equation*}
又因为与所有$\boldsymbol{x}\in\mathbb{R}^n$都正交的只有零向量，所以这样$\boldsymbol{a}$唯一.\qedhere
\end{proof}

\begin{ex}
  设 \(\boldsymbol{a}_1, \boldsymbol{a}_2, \boldsymbol{a}_3\)
  为 \(\mathbb{R}^3\) 的标准正交基.
  证明
  \begin{equation*}
    \boldsymbol{b}_1 = \frac{1}{3}(2\boldsymbol{a}_1 + 2\boldsymbol{a}_2 - \boldsymbol{a}_3), \quad
    \boldsymbol{b}_2 = \frac{1}{3}(2\boldsymbol{a}_1 - 1\boldsymbol{a}_2 + 2\boldsymbol{a}_3), \quad
    \boldsymbol{b}_1 = \frac{1}{3}(\boldsymbol{a}_1 - 2\boldsymbol{a}_2 - 2\boldsymbol{a}_3)
  \end{equation*}
  也是 \(\mathbb{R}^3\) 的标准正交基.
\end{ex}

\begin{proof}
  将条件中等式改写为
  \begin{equation*}
    [\boldsymbol{b}_1 \; \boldsymbol{b}_2 \; \boldsymbol{b}_3] =
    [\boldsymbol{a}_1 \; \boldsymbol{a}_2 \; \boldsymbol{a}_3] \cdot
    \frac{1}{3}
    \begin{bmatrix}
      2 & 2 & 1 \\
      2 &  -1 & -2 \\
      -1 & 2 & -2
    \end{bmatrix}
  \end{equation*}
  令
  \(
  Q =\frac{1}{3}
  \begin{bmatrix}
    2 & 2 & 1 \\
    2 &  -1 & -2 \\
    -1 & 2 & -2
  \end{bmatrix}
  \),
  以及 \(B = [\boldsymbol{b}_1 \; \boldsymbol{b}_2 \; \boldsymbol{b}_3]\).
  则断言等价于 \(B^{\mathrm{T}} \cdot B = \mathrm{I}_3\).
  由于 \(A = [\boldsymbol{a}_1 \; \boldsymbol{a}_2 \; \boldsymbol{a}_3]\)
  是正交方阵, 只需再证明 \(Q\) 是正交方阵即可. 直接计算即可.
\end{proof}

\begin{ex}
  向量 \(\boldsymbol{a}, \boldsymbol{b} \in \mathbb{R}^n\)
  围成一个三角形. 证明它的面积的平方为
  \(\frac{1}{4}(\|\boldsymbol{a}\|^{2} \cdot \|\boldsymbol{b}\|^2 - (\boldsymbol{a}^{\mathrm{T}}\boldsymbol{b})^2)\).
\end{ex}

\begin{proof}
  高向量为 \(\boldsymbol{a} - \frac{\boldsymbol{a}^{\mathrm{T}}\boldsymbol{b}}{\|\boldsymbol{b}\|^2}\boldsymbol{b}\).
  其长度平方 (即三角形关于边 \(\boldsymbol{b}\) 的高) 为
  \begin{align*}
    \left(\boldsymbol{a} - \frac{\boldsymbol{a}^{\mathrm{T}}\boldsymbol{b}}{\|\boldsymbol{b}\|^2}\boldsymbol{b}\right)^{\mathrm{T}}\left(\boldsymbol{a} - \frac{\boldsymbol{a}^{\mathrm{T}}\boldsymbol{b}}{\|\boldsymbol{b}\|^2}\boldsymbol{b}\right)
    & = \|\boldsymbol{a}\|^2 - 2(\boldsymbol{a}^{\mathrm{T}}\boldsymbol{b})^2 \cdot \|b\|^{-2}  + (\boldsymbol{a}^{\mathrm{T}}\boldsymbol{b})^2 \|\boldsymbol{b}\|^{-2}\\
    &= \|\boldsymbol{a}\|^2 - (\boldsymbol{a}^{\mathrm{T}}\boldsymbol{b})^2 \cdot \|\boldsymbol{b}\|^{-2}.
  \end{align*}
  因此四倍面积平方为 \(\|\boldsymbol{b}\|^2\) 乘以高的长度平方, 即为所求.
\end{proof}

\begin{ex}
  对向量组
  \begin{equation*}
    \boldsymbol{v}_1 =
    \begin{bmatrix}
      1 \\ 2 \\ -2
    \end{bmatrix},
    \boldsymbol{v}_2 =
    \begin{bmatrix}
      4 \\ 4 \\-3
    \end{bmatrix},
    \boldsymbol{v}_3 =
    \begin{bmatrix}
      -1 \\ 7 \\ 2
    \end{bmatrix}
  \end{equation*}
  进行 Schmidt 正交化.
\end{ex}

\begin{proof}[解]
  \begin{enumerate}
\item \(\widetilde{\boldsymbol{q}}_{1} = \boldsymbol{v}_{1}=\begin{bmatrix}1 \\ 2 \\-2\end{bmatrix}\).
\item \(\widetilde{\boldsymbol{q}}_{2} = \boldsymbol{v}_{2} - \mathrm{proj}_{\widetilde{\boldsymbol{q}}_{2}}(\boldsymbol{v}_{2})\)
这里 \(\mathrm{proj}_{\widetilde{\boldsymbol{q}}_{2}}\) 系 \(\boldsymbol{v}_{2}\)在 \(\widetilde{\boldsymbol{q}}_{2}\)
方向上的投影:
\begin{equation*}
  \mathrm{proj}_{\widetilde{\boldsymbol{q}}_1}(\boldsymbol{v}_2) =
  \frac{\boldsymbol{v}_2^{\mathrm{T}}\widetilde{\boldsymbol{q}}_1}{\widetilde{\boldsymbol{q}}_1^{\mathrm{T}}\widetilde{\boldsymbol{q}}_1}\widetilde{\boldsymbol{q}}_1.
\end{equation*}
计算可知, \(\widetilde{\boldsymbol{q}}_{2} = \begin{bmatrix}4\\4\\-3\end{bmatrix}-\frac{18}{9}\begin{bmatrix}1\\2\\-2\end{bmatrix}=\begin{bmatrix}2\\0\\1\end{bmatrix}\).
\item \(\widetilde{\boldsymbol{q}}_{3} = \boldsymbol{v}_{3} -\mathrm{proj}_{\widetilde{\boldsymbol{q}}_{1}}(\boldsymbol{v}_{3}) - \mathrm{proj}_{\widetilde{\boldsymbol{q}}_{2}}(\boldsymbol{v}_{3})\).
计算可知,
\begin{align*}
  \widetilde{\boldsymbol{q}}_3 &=
  \begin{bmatrix}
    -1 \\ 7 \\ 2
  \end{bmatrix}
  -\frac{9}{9}
  \begin{bmatrix}
    1 \\ 2 \\ -2
  \end{bmatrix}
  -\frac{0}{5}
  \begin{bmatrix}
    2 \\ 0 \\ 1
  \end{bmatrix} \\
  &=
  \begin{bmatrix}
    -2 \\ 5 \\ 4
  \end{bmatrix}.
\end{align*}
\end{enumerate}
最后, 归一化, 可得
\begin{equation*}
  [\boldsymbol{q}_1,\boldsymbol{q}_2,\boldsymbol{q}_3] =
  \begin{bmatrix}
    1/3 & 2/\sqrt{5} & -2/3\sqrt{5} \\
    2/3 & 0 & 5/3\sqrt{5}\\
    -2/3 & 1/\sqrt{5} & 4/3\sqrt{5}
  \end{bmatrix}.
\end{equation*}
\end{proof}

\begin{ex}
  求矩阵 \(\begin{bmatrix} 2 & 1 & -1 \\ 1 & 0 & 2 \\ 2 & -1 & 3 \end{bmatrix}\) 的 QR 分解.
\end{ex}

\begin{proof}[解]
  设矩阵的列分别为 \(\boldsymbol{a}_{1}, \boldsymbol{a}_{2}, \boldsymbol{a}_{3}\).
  使用 Gram-Schmidt 得
\begin{enumerate}
\item \(\widetilde{\boldsymbol{q}}_{1} = \boldsymbol{a}_{1} = \begin{bmatrix} 2 \\ 1 \\2 \end{bmatrix}\).
\item \(\widetilde{\boldsymbol{q}}_{2} = \boldsymbol{a}_2 - \mathrm{proj}_{\widetilde{\boldsymbol{q}}_1}(\boldsymbol{a}_2) = \boldsymbol{a}_2 - \frac{0}{9}\boldsymbol{a}_1 = \begin{bmatrix} 1 \\ 0 \\-1\end{bmatrix}\).
\item \(\widetilde{\boldsymbol{q}}_3 = \boldsymbol{a}_3 - \mathrm{proj}_{\widetilde{\boldsymbol{q}}_1}(\boldsymbol{a}_3) - \mathrm{proj}_{\widetilde{\boldsymbol{q}}_2}(\boldsymbol{a}_3)=\boldsymbol{a}_3-\frac{6}{9}\widetilde{\boldsymbol{q}}_1 - \frac{-4}{2} \widetilde{\boldsymbol{q}}_2 = \frac{1}{3}\begin{bmatrix}-1\\4\\-1\end{bmatrix}\).
\end{enumerate}
因此
\begin{align*}
  [\boldsymbol{a}_1 \; \boldsymbol{a}_2 \; \boldsymbol{a}_3]
  & = [\widetilde{\boldsymbol{q}}_1\; \widetilde{\boldsymbol{q}}_2 \; \widetilde{\boldsymbol{q}}_{3}]
  \begin{bmatrix}
    1 & 0 & 2/3 \\
    0 & 1 & -2 \\
    0 & 0 & 1
  \end{bmatrix} \\
  &=
  \begin{bmatrix}
    2/3 & 1/\sqrt{2} &  -1/3\sqrt{2} \\
    1/3 & 0          &  4/3\sqrt{2} \\
    2/3 & -1/\sqrt{2} &  -1/3\sqrt{2}
  \end{bmatrix}
  \cdot
  \begin{bmatrix}
    3 \\ {} & \sqrt{2} \\ {} & {} & \sqrt{2}
  \end{bmatrix}
  \cdot
  \begin{bmatrix}
    1 & 0 & 2/3 \\
    0 & 1 & -2 \\
    0 & 0 & 1
  \end{bmatrix} \\
  &=
    \begin{bmatrix}
      2/3 & 1/\sqrt{2} &  -1/3\sqrt{2} \\
      1/3 & 0          &  4/3\sqrt{2} \\
      2/3 & -1/\sqrt{2} &  -1/3\sqrt{2}
    \end{bmatrix}
    \cdot
    \begin{bmatrix}
      3 & 0 & 2 \\
      0 & \sqrt{2} & -2\sqrt{2} \\
      0 & 0 & \sqrt{2}
    \end{bmatrix} \\
  &=QR.
\end{align*}
\end{proof}
\end{document}
